\chapter{Le vie sensitive e motorie}

% ---------------------------------------------------------------------------
\section{La corteccia somatosensitiva: il lobo parietale}

Le aree corticali di riferimento per la sensibilità sono:
\begin{itemize}
  \item area sensitiva (somestesica) primaria (\textbf{S1});
  \item area sensitiva secondaria (\textbf{S2});
  \item area sensitiva supplementare;
  \item complesso amigdala, insula e parte anteriore del cingolo.
\end{itemize}
Le prime tre aree si trovano tutte nel lobo parietale.

Il \textbf{lobo parietale} è situato dietro la scissura di Rolando e sopra
la scissura silviana; anteriormente ha il lobo frontale, posteriormente
quello occipitale e inferiormente quello temporale, con i quali ha
rapporti anatomo-funzionali.

\rubrica{Localizzazione delle aree sensitive}
\begin{itemize}
  \item \textbf{area somestesica primaria} (S1): dietro il solco centrale
    $\rightarrow$ regione post-centrale;
  \item \textbf{area sensitiva secondaria} (S2): sul labbro superiore
    della scissura silviana;
  \item \textbf{area sensitiva supplementare}: sulla faccia mediale del
    lobo parietale;
  \item \textbf{complesso amigdala, insula e parte anteriore del
    cingolo}: coinvolto per la sola via paleo-spino-talamica.
\end{itemize}

% ---------------------------------------------------------------------------
\section{Organizzazione anatomo-funzionale delle vie sensitive}

Le vie sensitivo-sensoriali sono organizzate in modo da mantenere una
rappresentazione corporea ordinata per il loro intero decorso:
\begin{itemize}
  \item \textbf{segmentaria} (radici e midollo);
  \item \textbf{somatotopica} (talamo e corteccia): homuncoli sensitivi.
\end{itemize}

\rubrica{Organizzazione segmentaria}
Presente nelle radici dorsali e anche nel midollo spinale. Rispetto alle
radici dorsali, nel midollo l'organizzazione segmentaria è speculare per i
due sistemi sensitivo-sensoriali.

\rubrica{Organizzazione somatotopica}
Presente nel talamo e in corteccia (\textit{homunculus}):
\begin{enumerate}
  \item il somatotopismo è proporzionale alla densità di innervazione di
    ciascun segmento corporeo (non alla sua estensione metrica);
  \item la capacità di discriminazione sensitiva spaziale è proporzionale
    all'estensione delle rappresentazioni somatotopiche centrali.
\end{enumerate}

Ogni stazione sinaptica delle vie sensitive possiede, inoltre, circuiti
(soprattutto locali) di potenziamento del segnale, come pure circuiti (sia
segmentali che sovra-segmentali) di controllo dello stesso.

% ---------------------------------------------------------------------------
\section{Le vie sensitive}

Le vie sensitive principali sono tre:
\begin{enumerate}
  \item sistema delle colonne dorsali: tatto fine e discriminativo;
  \item fascio spinotalamico: dolore e temperatura;
  \item fascio spinocerebellare.
\end{enumerate}

\rubrica{Sistema delle colonne dorsali}
Trasmette informazioni dai meccanocettori e propriocettori alla corteccia
somatosensitiva, seguendo il percorso: fascicolo gracile / nucleo gracile
/ talamo, oppure fascicolo cuneato / nucleo cuneato / talamo. Decussa nel
bulbo.

\rubrica{Vie somatosensoriali}
Trasportano separatamente informazioni sul dolore, temperatura, tatto,
pressione, vibrazione, posizione e movimento degli arti. Sono formate da
tre neuroni:
\begin{itemize}
  \item \textbf{I ordine}: dai recettori sensoriali al midollo spinale o
    al tronco dell'encefalo;
  \item \textbf{II ordine}: dal midollo spinale al talamo, con
    decussazione delle fibre;
  \item \textbf{III ordine}: dal talamo all'area somato-sensoriale
    primaria ipsilaterale.
\end{itemize}

\rubrica{Neuroni sensitivi di 1°, 2° e 3° ordine}
\begin{description}
  \item[Neuroni sensitivi di 1° ordine:] originano in periferia; entrano
    nel midollo spinale attraverso le radici dorsali ed eventualmente il
    corno posteriore; entrano nel tronco encefalico per le sensazioni del
    viso.
  \item[Neuroni sensitivi di 2° ordine:] fanno sinapsi con quelli di 1°
    ordine nel corno posteriore del midollo spinale (tratto
    spino-talamico) oppure nel tronco (sistema delle colonne dorsali); il
    neurone di 2° ordine decussa prima di salire per raggiungere il
    talamo.
  \item[Neuroni sensitivi di 3° ordine:] fanno sinapsi con quelli di 2°
    ordine nel talamo e viaggiano fino a una specifica area della
    corteccia cerebrale.
\end{description}

\rubrica{Fascio spinotalamico}
Trasmette informazioni da meccanocettori, termocettori e nocicettori al
talamo; decussa nel midollo spinale.
\begin{itemize}
  \item \textbf{fascicolo spino-talamico laterale}: trasmette sensazioni
    di dolore e temperatura;
  \item \textbf{fascicolo spino-talamico anteriore}: trasmette sensazioni
    di tatto grossolano e pressione.
\end{itemize}

% ---------------------------------------------------------------------------
\section{Il sistema nervoso somatico e l'integrazione nervosa}

Il sistema nervoso regola le risposte fisiologiche agli eventi
perturbanti. Le azioni del sistema nervoso comprendono:
\begin{itemize}
  \item risposte somatiche agli stimoli esterni;
  \item integrazione sensori-motoria (sistema nervoso somatico);
  \item risposte viscerali agli stimoli interni;
  \item integrazione viscerale (sistema nervoso autonomo).
\end{itemize}

\rubrica{Elementi del sistema di integrazione somatico}
\begin{itemize}
  \item \textbf{recettori}: corpuscoli o terminazioni nervose che
    rilevano variazioni dell'ambiente interno;
  \item \textbf{vie afferenti}: gangli localizzati nei diversi punti del
    cranio (nervi sensoriali del sistema nervoso autonomo);
  \item \textbf{reti di elaborazione}: distribuite lungo tutto il
    nevrasse (numerosi nuclei localizzati nel tronco dell'encefalo, nella
    corteccia, nel sistema limbico e nel midollo spinale);
  \item \textbf{vie efferenti}: fibre dei motoneuroni, organizzate in un
    duplice ordine di neuroni che fanno sinapsi in un ganglio;
  \item \textbf{effettori}: fibre muscolari striate scheletriche (oppure
    fibre muscolari lisce, miocardiche o cellule ghiandolari).
\end{itemize}

% ---------------------------------------------------------------------------
\section{Il midollo spinale come livello di integrazione}

Il midollo spinale è il livello più periferico dell'integrazione: i
circuiti spinali che mediano le vie riflesse non coinvolgono i centri
superiori. La \textbf{sostanza grigia} contiene i corpi cellulari e i
prolungamenti dendritici e assonali, mentre la \textbf{sostanza bianca} è
formata dalle fibre mieliniche. Il midollo si suddivide in segmenti
chiamati \textit{mielomeri}, corrispondenti alle vertebre.

Le vie di ingresso sono costituite dalle radici dorsali: sono neuroni i
cui somi sono situati nei gangli delle radici dorsali. Le vie di uscita
sono costituite dalle radici ventrali: sono motoneuroni i cui somi sono
situati nelle corna anteriori (risposte somatiche) e nei nuclei laterali
(risposte viscerali).

% ---------------------------------------------------------------------------
\section{Le vie motorie}

\rubrica{Via cortico-spinale e fasci discendenti nel midollo spinale}
Assicurano il controllo cosciente volontario della muscolatura
scheletrica.
\begin{itemize}
  \item \textbf{fascio cortico-bulbare}: diretto ai nuclei motori dei
    nervi encefalici, assicura il controllo volontario di muscoli
    scheletrici oculari, masticatori e mimici, di alcuni muscoli del
    collo e della faringe;
  \item \textbf{fascio cortico-spinale laterale} e \textbf{fascio
    cortico-spinale anteriore}: diretti ai muscoli scheletrici tramite le
    piramidi bulbari, dove avviene la decussazione delle piramidi.
\end{itemize}

% ---------------------------------------------------------------------------
\section{Le vie extrapiramidali}

Queste vie possono modificare o indirizzare quadri motori somatici.

\rubrica{Funzioni dei principali fasci extrapiramidali}
\begin{itemize}
  \item \textbf{fasci vestibolo-spinali}: risposte motorie involontarie
    per il mantenimento della postura e dell'equilibrio;
  \item \textbf{fasci tetto-spinali}: regolazione della posizione di
    occhi, testa, collo e braccia in risposta a stimoli luminosi intensi,
    movimenti improvvisi o rumori forti;
  \item \textbf{fasci reticolo-spinali}: regolazione involontaria delle
    attività riflesse;
  \item \textbf{fasci rubro-spinali}: regolazione involontaria dei
    movimenti delle parti distali degli arti superiori, coadiuvando il
    controllo dei muscoli che eseguono movimenti più precisi.
\end{itemize}

I fasci vestibolo-spinale, tetto-spinale e reticolo-spinale coadiuvano
inoltre il controllo dei movimenti grossolani del tronco e dei muscoli
delle parti prossimali degli arti.

\rubrica{Nuclei encefalici coinvolti nelle vie extrapiramidali}
\begin{itemize}
  \item nuclei della base: nucleo caudato, putamen, globus pallidus;
  \item nucleo rosso;
  \item formazione reticolare;
  \item nucleo vestibolare;
  \item tetto (collicolo superiore e inferiore).
\end{itemize}

% ---------------------------------------------------------------------------
\section{Le sindromi extrapiramidali}

Importanti indicazioni derivano dall'osservazione dei fenomeni morbosi che
caratterizzano le sindromi extrapiramidali, nelle quali si rilevano
disordini del tono, della postura e del movimento.

\rubrica{Principali aspetti della patologia extrapiramidale}
\begin{itemize}
  \item aumento dei riflessi posturali;
  \item turbe del tono muscolare (rigidità, ipotonia, distonia);
  \item riduzione della motilità spontanea e perdita di movimenti
    automatici o associati;
  \item movimenti involontari anormali;
  \item forza conservata: i movimenti volontari, seppure perturbati,
    possono comunque essere realizzati.
\end{itemize}

\rubrica{Funzioni del sistema extrapiramidale}
\begin{itemize}
  \item regolazione del tono posturale;
  \item preparazione degli atteggiamenti tonici di predisposizione ai
    movimenti volontari;
  \item esecuzione di movimenti associati che rendono più agevoli e
    corretti i movimenti volontari;
  \item controllo delle modificazioni automatiche del tono e dei
    movimenti;
  \item riflessi che accompagnano le situazioni affettive e attentive;
  \item controllo dei movimenti originariamente volontari, poi divenuti
    automatici attraverso l'esercizio (es.\ marciare, scrivere);
  \item inibizione di movimenti involontari (ipercinesie), liberati
    invece nelle malattie extrapiramidali.
\end{itemize}
